\chapter{Implementation}
\label{ch:implementation}

This chapter documents how the five core features defined in Chapter~\ref{ch:analysis} were
translated into working code. Each feature section follows the same structure: the purpose and
technical challenge are contextualised first, the application screenshot shows the feature in
action, the underlying logic is explained in prose, and a targeted code listing illustrates the
most important implementation decision. Section~\ref{sec:impl-testing} covers automated and
usability testing, and section~\ref{sec:impl-project-org} describes the project organisation.

% ----------------------------------------------------------------
\section{Course Administration Module (FR-01)}
\label{sec:impl-course}

The course administration feature addresses FR-01: Instructors must be able to create and manage
courses, upload Course Materials into named sections, and control student enrollment. The
technical challenge in this feature is keeping the file storage backend transparent to the
business logic so that the development environment (local disk) and the production environment
(Google Cloud Storage) can be swapped by changing a Spring profile, without any changes to the
service layer.

% TODO: Insert screenshot Figure 4.1
% \begin{figure}[htbp]
%     \centering
%     \includegraphics[width=0.9\textwidth]{images/fig_4_1_course-admin.png}
%     \caption{Instructor course management page: section list with upload controls}
%     \label{fig:course-admin}
% \end{figure}

The course management page presents the Instructor with a list of sections and a file upload
button per section. Selecting a file triggers a \texttt{multipart/form-data} request to the
backend; the Instructor does not see storage paths or file types, which are resolved automatically
by the backend.

The \texttt{CourseService} delegates all file operations to the \texttt{FileStorageService}
interface. When \texttt{uploadMaterial} is called, the service stores the file, infers the
\texttt{MaterialType} from the file extension, and persists the resulting URL in the
\texttt{CourseMaterial} entity. The ownership check in \texttt{assertInstructor} ensures that
only the course owner can modify sections or materials; any other authenticated user receives an
HTTP~403 response. The listing below shows the upload method, which encapsulates this flow in
under twenty lines.

\begin{lstlisting}[language=Java, caption={Code listing 4.1 --- Course material upload with storage abstraction (CourseService.java)}, label=lst:course-upload]
public CourseMaterialDTO uploadMaterial(
        Long sectionId, Long requestingUserId,
        MultipartFile file, String title) {
    CourseSection section = findSection(sectionId);
    // Only the course instructor may upload
    assertInstructor(section.getCourse(), requestingUserId);

    // FileStorageService is injected; impl selected by Spring profile
    String fileUrl = fileStorageService.store(file, "courses");
    MaterialType fileType = detectFileType(
            file.getOriginalFilename(), file.getContentType());

    CourseMaterial material = CourseMaterial.builder()
            .title(title).fileUrl(fileUrl)
            .fileType(fileType)
            .originalFileName(file.getOriginalFilename())
            .section(section).build();
    return toMaterialDTO(materialRepository.save(material));
}
\end{lstlisting}

The \texttt{FileStorageService} interface defines three methods (\texttt{store}, \texttt{copy},
\texttt{delete}), and the active implementation is injected by Spring. This design satisfies
NFR-04 (maintainability) by keeping the service layer unaware of the storage backend, and means
that switching to a cloud provider for production requires only an environment variable change,
not a code change.

% ----------------------------------------------------------------
\section{Content Extension System (FR-02)}
\label{sec:impl-fork}

The Content Extension System addresses FR-02: students must be able to clone a course into a
Private Workspace, modify it freely, and submit a Merge Proposal for the Instructor to review.
The key technical challenges are the copy-on-write cloning strategy (avoiding redundant file
copies for shared base materials) and the referential-integrity constraint that prevents hard
deletion of materials that are still referenced by pending proposals.

% TODO: Insert screenshot Figure 4.2
% \begin{figure}[htbp]
%     \centering
%     \includegraphics[width=0.9\textwidth]{images/fig_4_2_content-extension.png}
%     \caption{Student Private Workspace view with Merge Proposal submission dialog}
%     \label{fig:content-extension}
% \end{figure}

The workspace view shows the student a mirror of the course structure. An upload button per
section allows the student to add new materials. The Merge Proposal dialog appears when the
student selects one or more materials and clicks the contribute button; the student picks a
target section (or names a new one) and submits.

When \texttt{forkCourse} is called on \texttt{WorkspaceService}, the backend iterates over all
sections and Course Materials of the source course and creates \texttt{WorkspaceSection} and
\texttt{WorkspaceMaterial} rows that point to the original file URLs (\texttt{isReference = true})
rather than copying the files. Only materials uploaded by the student later are stored as new
files (\texttt{isReference = false}). When a student removes a reference material from their view,
\texttt{deleteMaterial} in \texttt{WorkspaceService} sets \texttt{isHidden = true} on the row
instead of deleting it, because the row may still be referenced by a \texttt{ContributionProposal}
foreign key. The listing below shows this soft-delete logic.

\begin{lstlisting}[language=Java, caption={Code listing 4.2 --- Soft-delete for reference materials (WorkspaceService.java)}, label=lst:soft-delete]
public void deleteMaterial(Long materialId, Long userId) {
    WorkspaceMaterial material = materialRepository
            .findById(materialId)
            .orElseThrow(() ->
                new RuntimeException("Material not found: " + materialId));
    assertOwner(material.getSection().getSpace().getWorkspace(), userId);

    if (Boolean.TRUE.equals(material.getIsReference())) {
        // Soft-delete: keep the row so ContributionProposal FK stays intact
        material.setIsHidden(true);
        materialRepository.save(material);
    } else {
        // Hard-delete: student's own file, safe to remove
        fileStorageService.delete(material.getFileUrl());
        materialRepository.delete(material);
    }
}
\end{lstlisting}

The two-branch delete strategy preserves referential integrity without requiring a separate
audit table or a deferred constraint. The \texttt{WorkspaceSectionDTO} mapper already filters
hidden materials before returning the section content to the frontend, so the student's view is
consistent with the logical deletion even though the database row remains.

% ----------------------------------------------------------------
\section{Contextual Messaging System (FR-03)}
\label{sec:impl-chat}

The Contextual Messaging System addresses FR-03: users must be able to post messages in the
workspace chat and attach a Contextual Anchor to a specific Course Material so that other users
can open the material from within the message. The technical challenge is embedding a non-editable
interactive chip into a plain-text-like input field so that the anchor is preserved through
submission, serialisation, and re-render.

% TODO: Insert screenshot Figure 4.3
% \begin{figure}[htbp]
%     \centering
%     \includegraphics[width=0.9\textwidth]{images/fig_4_3_contextual-chat.png}
%     \caption{Workspace chat panel: message with Contextual Anchor chip and autocomplete dropdown}
%     \label{fig:contextual-chat}
% \end{figure}

The chat panel is shown after a student typed \texttt{@lecture} and the autocomplete dropdown
appeared. After selecting a Course Material, a styled chip replaces the \texttt{@}-query in the
input. The submitted message is rendered with a clickable badge that opens the file.

The message input uses a \texttt{<div contentEditable>} element rather than a
\texttt{<textarea>} because \texttt{<textarea>} can only hold plain text, whereas a
\texttt{contentEditable} div can host arbitrary DOM nodes. When the user selects a Course
Material from the autocomplete, the \texttt{insertTag} function deletes the \texttt{@<query>}
text and inserts a \texttt{<span>} element with \texttt{contentEditable="false"} and
\texttt{data-materialId}, \texttt{data-materialTitle}, and \texttt{data-materialUrl} attributes.
Before submission, \texttt{parseInput} serialises the span into the token
\texttt{@[<id>:<title>]}, which is stored as the message text. On render, a regular expression
parses the token back into a \texttt{<Badge>} component. The listing below shows the serialisation
step.

\begin{lstlisting}[language=Java, caption={Code listing 4.3 --- DOM serialisation to portable anchor token (contextual-chat.tsx)}, label=lst:anchor-token]
// Walks the contentEditable div's child nodes and builds the token string
function parseInput(div: HTMLDivElement): string {
  let result = "";
  div.childNodes.forEach((node) => {
    if (node.nodeType === Node.TEXT_NODE) {
      result += node.textContent ?? "";          // plain text
    } else if (node instanceof HTMLElement
               && node.dataset.materialId) {
      // Non-editable chip -> portable anchor token
      result += `@[${node.dataset.materialId}:${node.dataset.materialTitle}]`;
    }
  });
  return result.trim();
}
\end{lstlisting}

The token format \texttt{@[id:title]} is compact, survives JSON serialisation, and is
unambiguous to parse with a single regular expression on the render side. Storing the material
ID inside the token allows the backend to validate or re-fetch the material metadata if needed,
while the title provides a human-readable fallback if the material is later renamed.

\subsection{Workspace Collaboration}
\label{subsec:impl-workspace-collaboration}

To facilitate peer-to-peer discussion, Private Workspaces must support collaboration among multiple students. The system achieves this by allowing the workspace owner to add other users as members. The \texttt{WorkspaceService} manages this by persisting a list of \texttt{SpaceMemberDTO} objects associated with the space. Before a user can post or view messages in the chat, the \texttt{assertMember} security check verifies their membership. This transforms the Private Workspace from an isolated clone into a collaborative study environment, directly supporting the interactive requirements of FR-03.

\subsection{Real-Time Communication via WebSockets}
\label{subsec:impl-chat-sockets}

For the discussion to be truly interactive, messages must be delivered to all workspace members instantly without client-side polling. This is implemented using WebSockets with STOMP over SockJS (as justified in section~\ref{subsec:websocket-tech}). While earlier iterations used WebSockets exclusively for study sessions, the infrastructure was extended to power the Contextual Messaging System.

When a user submits a message, the \texttt{SpaceChatController} intercepts the HTTP \texttt{POST} request, saves the message to the database, and immediately broadcasts the message payload to the STOMP destination corresponding to the workspace ID.

\begin{lstlisting}[language=Java, caption={Code listing 4.4 --- WebSocket broadcast for workspace chat (SpaceChatController.java)}, label=lst:socket-broadcast]
@PostMapping("/{spaceId}/messages")
public ResponseEntity<MessageDTO> postMessage(
        @PathVariable Long spaceId,
        @RequestBody MessageRequest request) {
    MessageDTO savedMessage = chatService.saveMessage(spaceId, request);
    
    // Broadcast to space members over WebSocket channel
    messagingTemplate.convertAndSend(
        "/topic/space/" + spaceId + "/chat", 
        savedMessage
    );
    return ResponseEntity.ok(savedMessage);
}
\end{lstlisting}

On the frontend, the React application uses a WebSocket hook to subscribe to the \texttt{/topic/space/\{spaceId\}/chat} channel. Upon receiving a message event, the client appends the new message to the Redux state, triggering an immediate UI update for all connected members.

% ----------------------------------------------------------------
\section{AI-Based Content Querying (FR-04)}
\label{sec:impl-ai}

The AI querying feature addresses FR-04: users must be able to ask natural-language questions
about the content of the currently open workspace and receive answers grounded in that content.
The technical challenges are normalising three distinct document URL formats into a single byte
stream, constructing a structured prompt that keeps the model grounded in the document, and
handling API rate limits gracefully so the student always receives a useful response.

% TODO: Insert screenshot Figure 4.4
% \begin{figure}[htbp]
%     \centering
%     \includegraphics[width=0.85\textwidth]{images/fig_4_4_ai-assistant.png}
%     \caption{AI assistant panel: student question with grounded answer and typing animation}
%     \label{fig:ai-assistant}
% \end{figure}

The AI panel shows the student's question at the top and the LLM-generated answer below,
displayed with a typewriter animation. The panel is collapsible so the student can read the
Course Material alongside the answer.

When the user submits a question, the frontend sends the question and the selected
\texttt{documentUrl} to \texttt{POST /api/chat/query}. \texttt{DocumentService} normalises the
URL (handling HTTP/HTTPS remote URLs, \texttt{file://} URIs, and bare local paths) and extracts
the plain text using Apache PDFBox~\cite{ApachePDFBox}. \texttt{GeminiService} then constructs a
three-part prompt (system instruction, document context truncated to 10\,000 characters, and the
user's question) and calls the Gemini API via the Google GenAI Java SDK~\cite{GoogleGenAI}. The
listing below shows the prompt construction.

\begin{lstlisting}[language=Java, caption={Code listing 4.5 --- Structured prompt construction (GeminiService.java)}, label=lst:gemini-prompt]
String systemInstruction =
    "You are a teaching assistant. Prefer the provided document " +
    "context over your own knowledge when it is relevant.";

// Truncate to control cost and stay within the context window
String truncatedContext = docText.length() > 10_000
    ? docText.substring(0, 10_000) + "\n[...document truncated...]"
    : docText;

String prompt = String.format(
    "Document context:\n%s\n\nStudent question: %s",
    truncatedContext, question);

GenerateContentResponse response = client
    .models.generateContent(MODEL, prompt,
        GenerateContentConfig.builder()
            .systemInstruction(systemInstruction).build());
\end{lstlisting}

The system instruction instructs the model to prefer the document context, which is the core
of the Retrieval-Augmented Generation (RAG) pattern implemented here. Truncating the context at
10\,000 characters controls per-request cost while keeping the most relevant part of the document
available to the model. When the API returns a \texttt{429 RESOURCE\_EXHAUSTED} status, the
service parses the retry delay and surfaces it to the user as a human-readable message rather
than propagating a raw exception.

% ----------------------------------------------------------------
\section{Access Control (FR-05)}
\label{sec:impl-access}

The access control feature addresses FR-05: the system must distinguish Instructors from Students
at the API level and reject unauthorised requests before they reach the service layer. The
technical challenge is making the authentication filter stateless so that no server-side session
store is required, while still providing the full user identity (including role) to every
downstream service call.

% TODO: Insert screenshot Figure 4.5
% \begin{figure}[htbp]
%     \centering
%     \includegraphics[width=0.85\textwidth]{images/fig_4_5_access-control.png}
%     \caption{Role-differentiated navigation: student view (left) versus Instructor view (right)}
%     \label{fig:access-control}
% \end{figure}

The two navigation sidebars illustrate the role distinction. The student sees course browsing
and workspace links; the Instructor additionally sees the course management and Merge Proposal
review links. Both views are generated from the same component by reading the role field from
the JWT stored in the React context.

When a user authenticates, \texttt{AuthService} validates the credentials, looks up the
\texttt{User} entity, and calls \texttt{JwtUtil.generateToken} to produce a signed JWT embedding
the user's email as the subject. The token is signed with HMAC-SHA256 and expires after ten
hours. On every subsequent request, \texttt{JwtAuthenticationFilter} extracts the token from the
\texttt{Authorization: Bearer} header, validates the signature and expiry, loads the
\texttt{UserDetails} by email, and populates the Spring Security context. The listing below
shows the core of the filter.

\begin{lstlisting}[language=Java, caption={Code listing 4.6 --- JWT validation in the authentication filter (JwtAuthenticationFilter.java)}, label=lst:jwt-filter]
// Runs once per request; short-circuits if no token is present
String authHeader = request.getHeader("Authorization");
if (authHeader != null && authHeader.startsWith("Bearer ")) {
    jwt = authHeader.substring(7); // strip "Bearer " prefix
}
if (jwt == null) { filterChain.doFilter(request, response); return; }

userEmail = jwtUtil.extractUsername(jwt); // decode subject claim
if (userEmail != null
        && SecurityContextHolder.getContext().getAuthentication() == null) {
    UserDetails userDetails =
        userDetailsService.loadUserByUsername(userEmail);
    if (jwtUtil.isTokenValid(jwt, userDetails)) {
        // Populate security context so @PreAuthorize annotations work
        UsernamePasswordAuthenticationToken authToken =
            new UsernamePasswordAuthenticationToken(
                userDetails, null, userDetails.getAuthorities());
        SecurityContextHolder.getContext().setAuthentication(authToken);
    }
}
filterChain.doFilter(request, response);
\end{lstlisting}

The filter populates the Spring Security context with the user's authorities (which encode the
role), enabling \texttt{@PreAuthorize} annotations and explicit \texttt{assertInstructor} checks
in the service layer to reject Student requests on Instructor-only operations. The stateless
approach means that the server holds no session state between requests, satisfying NFR-02 and
simplifying horizontal scaling.

% ----------------------------------------------------------------
\section{Testing}
\label{sec:impl-testing}

\subsection{Backend Automated Tests}
\label{subsec:backend-tests}

The backend uses three levels of automated tests, all implemented with JUnit~5 and the Spring
Boot Test framework.

\emph{Unit tests} cover individual service classes in isolation. Dependencies are substituted
with Mockito mocks so that each test exercises exactly one unit. Examples include
\texttt{UserServiceTest}, \texttt{StudyGroupServiceTest}, and \texttt{GamificationServiceTest}.

\emph{Repository tests} verify that the Spring Data JPA repositories produce correct SQL and
return expected entity graphs. They run against an in-process H2 database seeded with test data.
\texttt{UserRepositoryTest} checks that user look-up by email and username behaves correctly.

\emph{Integration tests} exercise complete request-response cycles through the Spring MVC layer.
\texttt{UserControllerIntegrationTest} and \texttt{StudySessionControllerIntegrationTest} spin
up a \texttt{@SpringBootTest} application context, issue HTTP requests via \texttt{MockMvc}, and
assert on the HTTP status code and the JSON response body. \texttt{GlobalExceptionHandlerTest}
verifies that the centralised exception handler maps domain exceptions to the correct HTTP status
codes.

During development, several important bugs were found and corrected through testing. A
referential-integrity violation when deleting reference materials was resolved by introducing the
soft-delete \texttt{isHidden} flag described in section~\ref{sec:impl-fork}. A scroll-hijacking
issue in the AI chat panel was discovered during manual end-to-end testing and fixed by
introducing a \texttt{userScrolledRef} guard that suppresses programmatic scroll when the user
has manually scrolled upward.

\subsection{Frontend Automated Tests}
\label{subsec:frontend-tests}

Complementing the backend testing suite, the frontend employs a dual-framework testing strategy to verify user interfaces and application flows.

\emph{Unit and integration tests} are executed using Vitest and the React Testing Library. Components such as the authentication forms are mounted in an isolated \texttt{jsdom} environment. By mocking the Next.js router and Redux store, the tests simulate user input and assert on the resulting DOM updates and validation errors without relying on a running server.

\emph{End-to-end (E2E) tests} are orchestrated with Playwright. These tests spin up a headless browser (such as Chromium) and navigate through complete user journeys. By leveraging Playwright's network interception capabilities, the tests mock the backend API responses. This approach ensures that the frontend behaves correctly under various backend states---such as a successful login or an HTTP 401 Unauthorized response---in a highly deterministic and fast-executing manner.

\subsection{User and Usability Tests}
\label{subsec:user-tests}

Each test below is written as a numbered step sequence. Pre-conditions and edge cases are noted
inline.

\bigskip
\noindent\textbf{Test 4.1 --- Upload and browse a Course Material (FR-01)}
\begin{enumerate}
    \item Log in as an Instructor account.
    \item Navigate to an existing published course and open the course management view.
    \item Select a section and click the upload button. Choose a PDF file of at least 1~MB.
    \item Confirm the upload. Expected: the material appears in the section list with the correct
          file name and a PDF icon.
    \item Log out. Log in as an enrolled Student. Navigate to the same course and section.
    \item Confirm the material is visible and can be downloaded.
\end{enumerate}

\noindent\textbf{Edge case}: Attempt the upload step (step 3) while logged in as a Student.
Expected: the upload button is absent from the UI; a direct API call to
\texttt{POST /api/courses/sections/{id}/materials} returns HTTP~403 Forbidden.

\bigskip
\noindent\textbf{Test 4.2 --- Clone course and submit Merge Proposal (FR-02)}
\begin{enumerate}
    \item Log in as a Student enrolled in a published course.
    \item Navigate to the course and click the ``Clone to workspace'' button.
    \item Confirm the Private Workspace is created with sections and materials mirroring the
          original course.
    \item Upload a new PDF into one of the workspace sections.
    \item Select the uploaded material, click ``Contribute'', choose an existing target section,
          optionally add a message, and submit the Merge Proposal.
    \item Log out and log in as the course Instructor. Navigate to the Merge Proposal review tab.
    \item Confirm the proposal appears with status \texttt{PENDING}. Approve it.
    \item Navigate to the course section. Confirm the approved material appears with the
          contributor's name.
\end{enumerate}

\noindent\textbf{Edge case}: In step~5, select a reference material (one cloned from the course)
rather than a newly uploaded file. Expected: the proposal is accepted, but the backend copies the
file from the original URL into the course storage, not from a non-existent student upload.

\bigskip
\noindent\textbf{Test 4.3 --- Post a message with Contextual Anchor (FR-03)}
\begin{enumerate}
    \item Log in as a Student and open a workspace.
    \item In the chat panel, type \texttt{@} followed by part of a Course Material name.
    \item Confirm the autocomplete dropdown appears with matching materials.
    \item Select a material. Confirm the chip is inserted inline in the message input.
    \item Submit the message. Confirm the chip is rendered as a clickable badge in the chat.
    \item Click the badge. Confirm the Course Material file is downloaded or opened.
\end{enumerate}

\noindent\textbf{Edge case}: Open the same workspace as a second enrolled Student. Confirm the
message with the Contextual Anchor is visible to the second student and the badge is functional.

\bigskip
\noindent\textbf{Test 4.4 --- Query AI assistant (FR-04)}
\begin{enumerate}
    \item Log in as a Student, open a workspace, and select a Course Material (PDF).
    \item Open the AI assistant panel and type a factual question about the document's content.
    \item Submit. Confirm the response appears with a typewriter animation and references
          concepts from the document rather than generic knowledge.
    \item Type a follow-up question. Confirm the AI response reflects the conversation history.
\end{enumerate}

\noindent\textbf{Edge case}: Open the AI panel without selecting any Course Material and submit
a question. Expected: the backend sends the question without document context and the AI responds
with general knowledge; no error is shown.

\bigskip
\noindent\textbf{Test 4.5 --- Role-based access control (FR-05)}
\begin{enumerate}
    \item Log in as a Student. Open the browser developer tools and copy the JWT from the
          \texttt{Authorization} header of any API request.
    \item Using \texttt{curl} or Swagger UI, call
          \texttt{POST /api/courses} (Instructor-only endpoint) with the Student JWT.
    \item Expected: the server returns HTTP~403 Forbidden with an error body.
    \item Repeat for \texttt{PUT /api/courses/\{id\}/sections}
          and \texttt{POST /api/proposals/\{id\}/review}.

    \item In each case confirm the response is 403 and the operation has no effect.
\end{enumerate}

\noindent\textbf{Edge case}: Submit a request with a malformed or expired JWT to any protected
endpoint. Expected: the server returns HTTP~401 Unauthorized with the message ``Unauthorized''.

% ----------------------------------------------------------------
\section{Project Organisation}
\label{sec:impl-project-org}

\subsection{Workflow and Version Management}
\label{subsec:workflow-vcs}

The project is maintained in a single Git repository with a monorepo layout separating the
\texttt{backend/}, \texttt{frontend/}, and \texttt{documentation/} directories. Development
follows a branch-based workflow: each feature or fix is developed on a dedicated branch and
merged into \texttt{main} via a pull request after passing all automated checks. Commits follow
the Conventional Commits specification to keep the history readable and to support automated
changelog generation.

\subsection{CI/CD Pipeline}
\label{subsec:cicd}

A GitHub Actions workflow runs on every push to the \texttt{main} branch and on every pull
request targeting it. The pipeline consists of three sequential jobs. The first job compiles the
backend with Maven and runs the full JUnit test suite; if any test fails, the remaining jobs are
skipped. The second job installs frontend dependencies with \texttt{npm ci} and runs the TypeScript
type-checker and linter. The third job builds both Docker images and verifies that the
\texttt{docker compose up} command completes without errors.

% TODO: Insert CI/CD pipeline diagram as Figure 4.6
% \begin{figure}[htbp]
%     \centering
%     \includegraphics[width=0.9\textwidth]{images/fig_4_6_cicd-pipeline.png}
%     \caption{GitHub Actions CI/CD pipeline: three-stage build, test, and integration check}
%     \label{fig:cicd-pipeline}
% \end{figure}

The pipeline enforces that \texttt{main} is always in a deployable state. No code reaches the
branch without passing the backend tests and the frontend type check, which together cover the
most common categories of regression.

\subsection{Documentation}
\label{subsec:documentation}

The thesis source is maintained as a LaTeX project using the FIT CTU thesis template. The
\texttt{ctufit-thesis.tex} root file includes the chapter files from the \texttt{text/}
subdirectory. Figures are stored under \texttt{documentation/images/} and referenced from the
chapter files. The bibliography is maintained in \texttt{text/bib-database.bib} and compiled
with Biber.

The backend API is documented automatically by Springdoc OpenAPI, which generates a Swagger UI
from the controller annotations and makes it available at \texttt{/swagger-ui/index.html} when
the application is running. Public service methods carry JavaDoc comments that explain the
purpose and contract of each operation. The \texttt{/docs/} directory in the repository root
serves as a staging area for feature-level implementation notes used during thesis writing.
